\chapter{客观要件一：行为}

% Source PDF physical page 42; see audit/source_coverage.json.
\begin{EditionOpeningPanel}
重要考点：被害人自陷风险，不作为犯中的“作为义务”。
\end{EditionOpeningPanel}

\section{危害行为}

[问题］狗剩明知卖淫女小美有严重性病，仍情不自禁，奋不顾身地飞蛾扑火，与其发生性行为，感染了严重性病（鉴定为轻伤）。小美是否需要对狗剩的轻伤负责？（求仁得仁案）

\subsection*{一、特征}

犯罪是一种危害行为。危害行为具有以下三项特征：

1.有体性。行为是人的身体活动，包括积极举动和消极静止。这是行为的客观特征。基于此，人的内心思想再邪恶，只要没有付诸行动，不是危害行为。思想无罪。

2.有意性。行为是人有意识实施的。这是行为的主观特征。基于此，梦游、癫痫、无意识的反射举动等不是刑法上的行为。

3.有害性。行为具有法益侵害性。这是行为的实质特征。基于此，现代刑法认为，同性恋等行为没有法益侵害性，不是刑法上的危害行为。

\subsection*{二、判断}

\subsubsection*{（一）危害行为与生活行为的区分}

\CJKunderwave{危害行为：行为对法益制造了危险}。生活行为：行为对法益没有制造危险。（2025年试题)

[注意】根据客观主义立场，定罪应先看客观上有无危害行为。行为人虽然主观上有犯意，但客观行为对法益没有制造危险的，不属于危害行为，即使偶然发生危害结果，也不能因此将其行为认定为危害行为。

[提示】考题用主观来诱导考生先人为主，是常考套路。

例1（2006年第13题），甲想杀死乙，心想：劝乙到大街上跑步，如果被车撞死才好。乙答应，跑步时竟真的被车撞死。甲的行为对乙的生命没有创设危险，不是危害行为。

例2，甲送乙一袋大枣，希望乙噎死。乙果真噎死。甲的行为不是危害行为。

% Source PDF physical page 43; see audit/source_coverage.json.
\subsubsection*{（二）降低危险与替代危险}

1.降低危险的行为不是危害行为。

例1，在山脚下，甲看到一块石头要落下来砸向乙的头部，甲推乙一把，石头砸中乙的肩膀。甲的行为降低了危险，不是危害行为。如果甲原本可以大力推一把，让石头不要砸中乙的肩膀，甲却故意想让石头砸中肩膀，也即甲有伤害故意，即便如此，甲也无罪，因为甲的行为不是危害行为。判断有罪无罪应先看客观行为。

例2，甲开车不慎撞伤丙，丙倒在马路中央，甲逃逸。乙路过，将丙从马路中央移至马路边上，然后离去。乙的行为不是危害行为。

2.替代危险。降低危险是降低原有危险的危害程度。替代危险，是指开创新的危险，不过新的危险比原有危险的危害程度低。

例如（2018年试题），邻居房屋发生火灾，甲冲进去欲救一个婴儿，但是由于无法走出去，只好将婴儿从二楼的窗户扔出去，导致婴儿重伤。甲开创了新的危险（伤害），但该危险比原有危险（烧死）程度低。由于甲开创了新的危险，因此其行为属于危害行为，具有法益侵害性，到了排除违法事由阶段，可根据紧急避险或推定的被害人承诺排除违法性，进而无罪。\footnote{降低危险与替代危险的区别，是德国刑法学的典型知识点，德国刑法教材有专门论述。我国传统刑法学没有该知识点。在我国刑法学知识实现现代化转型的大背景下，这些大陆法系典型知识点，近年来越来越多地被法考考查，是值得关注的现象。例如，违法性认识可能性、期待可能性等，均是我国传统刑法学没有的知识点，但法考均予以考查。不过，在借鉴德国刑法学理论时，必须考虑我国国情，不可照搬照抄。}

\subsubsection*{(三）被害人自陷风险（五星级考点)\protect\footnotemark}\footnotetext{参见张明楷：《刑法学中危险接受的法理》，载《法学研究》2012年第5期，第180页。}

被害人自陷风险的问题，涉及危害行为的判断、因果关系的判断。本书将该问题放在行为这讲阐述。一个实害结果的发生，有行为人与被害人的共同参与，该实害结果应算到谁的头上，由谁负责？判断的第一步是，判断危险的实行者、支配者是谁。

1.被害人是危险的实行者、支配者，行为人教唆或帮助被害人自陷风险。

\CJKunderwave{判断标准}：如果被害人满足以下两个条件，则被害人对结果负责（反之，由行为人负责)：

（1）主观上，被害人对危险有\CJKunderwave{认识能力}。如果被害人是未成年人或精神病患者，则认为其无认识能力。

（2）客观上，被害人对危险有\CJKunderwave{控制能力}，也即基于自己的意志自由，对危险具有控制、消除或避免的能力。

例1，甲为了逗乙，向湖中扔300元，乙跳入湖中去捡，被湖水淹死。乙的行为属于被害人自陷风险，与甲无关。

例2，丈夫赌博后回家，妻子不开门，丈夫在门外站一夜，欲用诚意打动妻子，被冻死。丈夫属于被害人自陷风险，自己负责。

例3，甲放火烧乙家房屋，在屋外的乙奋不顾身进去救自已的孩子，被烧死。乙对危险没有自由的控制能力。乙的死亡不应由乙负责，而应由甲负责。如果邻居、消防员等为了救乙的孩子，进去被烧死，也应由甲负责。（2011年第3题、2013年第52题）如果乙

% Source PDF physical page 44; see audit/source_coverage.json.
为了抢救屋里5000元，奋不顾身进去，被烧死，则多数说认为应由乙负责，少数说认为应由甲负责。（2012年主观题，2021年试题）

[结论】多数说认为，不看救援者的身份，而看被救对象的性质：救人或贵重物品，放火者负责；救普通物品，救援者自己负责。

[巩固练习]题1，甲女欲伤害男友乙，送其一双旱冰鞋，企盼其运动时摔伤。乙果真摔成重伤。乙的重伤由乙负责。（2013年第5题）

题2，甲想杀死乙，心想：乙是吸毒的，给乙赠送大量毒品，乙吸食过量死亡才好。甲在乙头脑清醒时向乙赠送了大量毒品，乙毒瘾发作时吸食过量死亡。甲的行为不是危害行为。如果甲在乙毒瘾发作、无控制能力时赠送毒品，乙吸食死亡，则甲的行为是危害行为，构成故意杀人罪。如果乙是未成年人或精神病患者，则甲的行为也是危害行为，对死亡结果应负责。（2007年第16题，2017年第52题）

题3，在本节开头“求仁得仁案”中，狗剩飞蛾扑火，属于被害人自陷风险，自己对轻伤负责。小美不需要对其轻伤负责，不构成故意伤害罪。但是，小美构成传播性病罪。

题4，甲女被派出所抓去。甲告知丈夫乙：给我弄瓶农药，我要喝点，吓唬警察。乙照办。甲喝后死亡。甲的死亡属于被害人自陷风险，乙无罪。

2.行为人是危险的实行者、支配者，被害人同意行为人的危险行为。

\CJKunderwave{判断标准}：如果行为人满足以下两个条件，行为人对结果负责（反之，由被害人负责）：

（1）主观上，行为人对危险有\CJKunderwave{认识能力}。

（2）客观上，行为人对危险有\CJKunderwave{控制能力}，也即基于自己的意志自由，对危险具有控制、消除或避免的能力。

[提示]虽然被害人同意行为人的危险行为，但同意接受危险行为不等于同意接受实害结果。况且，即使同意接受实害结果中的死亡结果或重伤结果，这种同意或承诺也是无效的。

例1，乙女因违法行为被警察带到派出所三楼，丈夫甲为了解救乙，经乙同意，将乙绑在绳子上，从厕所窗户外向下吊，中途绳子断裂，乙摔死。虽然有乙的同意，但甲的行为仍是危害行为，构成过失致人死亡罪。

例2，甲开车载着乙，看到前方湖面结冰，甲对乙说：我观察了，冰层挺厚，我们从冰面上开过去。乙答应。途中冰面破裂，乙被淹死。甲应对乙的死亡负责。

[巩固练习］题1，甲女让醉酒的男友乙开车送自已回家，乙不愿意，甲生气。乙为了讨甲开心，便醉酒驾车，由此发生车祸，导致甲死亡。虽然甲属于“不作不死”，但甲对乙的行为并没有实质的支配力，死亡不能归因于甲，而应归因于乙。甲乙构成危险驾驶罪的共同犯罪，乙是实行犯，甲是教唆犯。乙同时触犯交通肇事罪，择一重罪论处。

题2，在暴风雨中，乘客甲欲让摆渡工乙把自己渡过河。乙劝阻甲，指出此时渡河危险性很大。但甲因为有急事要办，执意要过河。乙只好冒险摆渡。船在河中被风浪掀翻。甲溺死，乙获救。甲仅是乘客，对乙不会形成支配力，因此，甲的死亡应由乙负责，乙构成过失致人死亡罪。

题3，甲醉酒骑摩托车，被警察乙拦住。乙命令甲骑摩托车载着自己去交警大队接受处理。甲只好照办，因醉驾，途中发生车祸，车祸导致后座的乙死亡。虽然甲是危险的实行者，但甲失去独立的意志自由，对危险没有控制能力，死亡结果不应由甲负责，而应由乙负责。

% Source PDF physical page 45; see audit/source_coverage.json.
\InsertDiagram{figures/manual_p045_1.tex}

\section{不作为犯的分类}

\InsertDiagram{figures/manual_p045_2.tex}

\subsection*{一、行为的分类}

1.作为，是指违反刑法禁止性规定的行为。例如，飞车抢夺的行为违反禁止抢夺他人财物的规定，就是作为。

2.不作为，是指违反刑法命令性规定的行为。例如，母亲不给婴儿喂养饿死婴儿的行为，违反抚养婴儿的命令性义务，属于不作为。

\subsection*{二、不作为犯的分类}

1.真正不作为犯

这是指刑法明文规定只能由不作为构成的犯罪。常考罪名：（1）丢失枪支不报罪（第129条）。（2）遗弃罪（第261条）。（3）拒不支付劳动报酬罪（第276条之一）。

（4）拒绝提供间谍犯罪、恐怖主义犯罪、极端主义犯罪证据罪（第311条）。（5）拒不执行判决、裁定罪（第313条）。（6）不解救被拐卖、绑架妇女、儿童罪（第416条）。这些罪名的共同特征是，立法者都设立了法定的作为义务而行为人不履行。

[注意]真正不作为犯中，不履行作为义务的方式，既可以是积极举动，也可以是消极静止。例如，不履行扶养义务的方式既可以是将老父亲扔到大街上，也可以是对躺在病床上的老父亲不闻不问。不能因为扔到大街上是积极举动，就认为遗弃罪是既可由作为构成也可由不作为构成的不真正不作为犯。

2.不真正不作为犯

这是指既可由作为构成也可由不作为构成的犯罪，当由不作为构成时，称之为不真正不作为犯。例如，故意杀人罪，既可用刀捅死人，也可将婴儿活活饿死。当用不作为饿死婴儿时称为不真正不作为犯。

（1）作为与不作为的判断

如何判断一个犯罪是由作为构成，还是由不作为构成？例如，甲开车不慎撞伤乙，

% Source PDF physical page 46; see audit/source_coverage.json.
①甲不救助乙，开车逃跑，乙因无人救助而死亡。②甲将乙藏在草丛中，然后离去，乙因无人救助而死亡。

判断方法：第一，不作为。这是指消极地不消除危险，也即有消除危险的义务，却消极地不履行该义务。依此，上述事例①是不作为，甲有救助义务，却不履行。

第二，作为。“作为”的第一个条件是，积极地制造危险。第二个条件是，制造的危险达到通常性危险的程度。通常性危险，是指通常情况下能够直接导致实害结果的危险。

例如，拿刀捅人心脏，通常情况下有致命危险，所以是作为方式的杀人。又如，甲将乙扔进河里，具有通常性的致命危险，所以是作为方式的杀人。而事例②，“藏在草丛中”虽然是个积极举动，但其本身不具有通常性的致命危险，所以本身不是作为方式的杀人。此时，如果甲又赶紧救乙，则不构成犯罪。甲不救乙而离去，属于不作为方式的杀人。

（2）作为与不作为可以竞合（2024年试题）

例如，主治医生甲故意撤除病人的救助装置，导致病人死亡。从救助的角度看，甲是不作为。从撤除装置的角度看，甲是作为。作为与不作为竞合，对此最终以作为论处。

3.持有型犯罪

[结论］持有型犯罪属于作为犯罪。例如，非法持有毒品罪、持有假币罪等都是作为犯罪。这是因为，其一，维持持有状态，不是消极无为，而需积极举动。其二，其行为本身就在直接侵害法益。

由于持有型犯罪不是不作为犯罪，不存在作为义务（上缴义务），因此不能认为，行为人没有上缴才构成犯罪。只要维持非法持有状态，就构成持有型犯罪。只要不维持非法持有状态，就不构成持有型犯罪。

例1（2006年第4题），甲下班回家后，发现自家门前放着一包来历不明、类似面粉的东西。甲第二天上班时拿到实验室化验，发现是海洛因，于是立即倒入厕所马桶冲入下水道。甲虽然没有将毒品上缴公安部门，但没有维持非法持有状态，故不构成非法持有毒品罪。

例2（2014年第57题），甲父去世前告诉甲“咱家院墙内埋着5支枪”，甲说“知道了”，但此后甲什么也没做，没有动这些枪。甲对枪支维持非法持有状态，构成非法持有枪支罪。

4.作为、不作为与故意、过失的关系

[结论］作为犯有故意犯罪，也有过失犯罪。不作为犯有故意犯罪，也有过失犯罪。不能认为：作为犯一定是故意犯罪，不作为犯一定是过失犯罪。

例如，警察疏忽大意，忘记履行职责，造成他人死亡，属于过失的不作为，构成玩忽职守罪（2019年试题）。警察看到甲在杀人，故意不阻止，属于故意的不作为犯，构成滥用职权罪。

\section{不作为犯的成立条件}

[问题］女友和母亲同时掉进河里，狗蛋不救母亲，而救女友小芳。母亲因无人施救而死亡。对狗蛋如何处理？（2015年第52题，考验男人案）

成立条件：（1）负有作为义务（应为）；（2）具有作为可能性（能为）；（3）具有结果避免可能性（不为）。

% Source PDF physical page 47; see audit/source_coverage.json.
\subsection*{一、负有作为义务 (应为)}

刑法上的作为义务是指，某个危险源对法益对象产生了危险流，刑法要求行为人去阻断危险流、保护法益对象。凭什么让行为人承担这种作为义务？其实质的理由根据有两种：一是对危险源的监管义务；二是对法益对象的保护义务。这种学说被称为\CJKunderwave{实质的二分}\CJKunderwave{说}，是法考的理论立场。

[提示】传统理论关于作为义务来源采取“形式的四分说”。该说过于简单和形式化理解，不足以应对法考真题考查。

\subsubsection*{（一）对危险源的监管义务}

某个危险源制造了危险，而行为人对危险源负有监管义务。

1.对危险物的监管义务。这里的危险物包括危险动物、危险物品、危险设施等。

例1，主人对饲养的凶狗负有监督义务，动物园的管理人对饲养的动物有管理义务。当动物咬人时，负有阻止义务。（2011年第52题，2024年试题）

例2，道路设施、电力设施、矿井、广告牌等负责人对这些设施、设备负有管理义务。

例3，机动车的所有人对机动车的使用负有管理义务。如果无驾照者、醉酒者、小孩等欲驾驶，则负有阻止义务。

例4，房东发现房屋有危险，有告知租客的义务。

2.对他人危险行为的监管义务。这里的他人与行为人具有监护、监管关系。

例1，父母对年幼子女的危险行为负有监督义务。如果年幼子女伤害别人，父母有阻止和救助义务。相反，儿子看到父亲在盗窃，儿子没有阻止义务。（2011年第52题，2014年第5题，2020年试题，2024年试题）

例2，家属对患狂躁症的家庭成员的危险行为负有监督义务。如果该狂躁症患者伤害别人，家属有阻止和救助义务。

例3，军官对士兵的危险行为有监管义务。

例4，幼儿园阿姨对小朋友的危险“恶作剧”负有监督义务。

例5，成年兄妹之间、夫妻之间没有监管关系。妻子对丈夫（税务局局长）的滥用职权、贪污受贿等行为没有阻止的义务。如果不阻止，不构成这些罪的不作为的帮助犯。（2013年第51题）

3.对自己先行行为的监管义务。自已的先行行为要产生作为义务，条件是：对法益创设了危险。

例1，某饭店的食物导致客人中毒，饭店管理人对客人有救助义务。

例2，汽车制造商、手机制造商、药品制造商对售出的有安全隐患的产品有召回义务。（2021年试题）

例3，甲在黑夜里将车停在高速路上，不采取措施以防止后面车辆“追尾”，导致车辆相撞，后车司机受伤。甲对受伤司机有救助义务。

例4，甲明知乙有癫痫，故意通过吵架引起乙癫痫发作，乙浑身抽搐。甲故意不救助，乙死亡。甲构成不作为的故意杀人罪。（2019年试题）

[巩固练习］题1，甲夜间在办公室用电热炉煮面条，不慎将文件碰到电热炉上燃烧，本可立即扑灭，却不扑灭，最后酿成重大火灾。甲构成不作为的放火罪。（2021年试题）

% Source PDF physical page 48; see audit/source_coverage.json.
题2，甲乙共同抢劫丙女，致丙昏迷。乙又欲强奸丙，甲负有阻止义务。甲不阻止，会构成不作为的帮助犯。

题3，大街上，男朋友向女朋友提出分手，女朋友声称如分手就割腕自杀。男朋友不制止，女朋友自杀身亡。由于提出分手不会给女朋友的生命创设危险，所以男朋友没有救助义务。当然，在道德情感上男朋友有救助义务，但这种道德情感义务在此不能上升为刑法上的义务。（2012年第4题）

题4，甲递给乙一把刀，让乙观赏。乙观赏时，忽然用来刺伤丙。甲对丙没有救助义务。

（1）正当防卫

第一，防卫人对不法侵害人造成伤害，对其有无救助义务？

例如（2013年第7题），甲对乙仅实施一般性伤害，乙防卫反击，致甲重伤倒地（此时不算防卫过当）。乙故意不予救助，甲死亡。

少数说：正当防卫行为不会产生救助义务。理由：先行行为必须是违法行为，才能产生作为义务。正当防卫行为是合法行为。合法行为不会产生作为义务。因此，乙没有救助义务。

多数说：防卫行为若会导致过当结果，则防卫人有防止过当结果发生的义务，否则构成防卫过当，应负刑事责任。先行行为不要求是违法行为。合法行为也能产生作为义务。本题中，如果甲死亡，则乙构成防卫过当。因此，乙有防止过当结果（甲死亡）发生的义务。乙故意不作为，导致死亡，构成不作为的故意杀人罪。

例如，甲欲杀乙，用刀砍乙，乙反击致甲重伤倒地。乙不予救助，甲死亡。由于甲砍杀乙，乙的防卫行为即使导致甲死亡，也不构成防卫过当。因此，乙没有救助甲的义务。乙的防卫行为构成正当防卫，不救助行为不构成不作为犯。

第二，防卫行为对第三方造成危险，谁有作为义务？多数说：不法侵害人和防卫人都有作为义务。少数说：只有不法侵害人有作为义务。

例如，甲故意伤害乙，乙为了防卫，向甲扔了一个火把，火把刚好掉入稻草堆里。乙趁甲转移注意力时立马逃离。甲见起火，不灭火，离开现场。火势变大，烧毁邻居家。按照多数说，甲有灭火义务，不灭火，构成不作为犯罪。乙在客观上制造了危险，也有灭火义务，但是，乙缺乏灭火的条件，缺乏作为可能性，因为甲在伤害乙，因此乙最终不构成不作为犯罪。（2020年、2023年试题）

（2）紧急避险可以成为先行行为而产生作为义务。例如，甲驾驶客运大巴，刹车失灵，为了避免车辆掉下悬崖，不得已撞伤路边的老大爷。甲的紧急避险行为虽然属于合法行为，但是给老大爷创设了危险，因此有救助义务。

（3）犯罪行为可以成为先行行为而产生作为义务。例如，甲盗伐林木，树木倒下砸晕乙，甲不救助而离去，后乙死亡。甲构成盗伐林木罪和不作为的故意杀人罪，数罪并罚。（2010年第52题）

\subsubsection*{（二）对法益对象的保护义务}

1.\CJKunderwave{特定关系}

基于特定关系，某项法益的保护依赖于行为人，当该法益处于危险境地时，行为人负有保护义务。这种特定关系主要由以下根据产生：

% Source PDF physical page 49; see audit/source_coverage.json.
（1）法律产生的特定关系

例1，路人甲看到附近大楼发生火灾，没有报警。虽然《消防法》规定任何人发现火灾都有报警义务，但是，报警义务不等于保护义务。无论根据何种法律，甲对大楼均没有保护义务，因此没有刑法上的作为义务。

例2，儿子落水，父亲有救助义务。在本节开头“考验男人案”中，母亲落水，子女有救助义务。女友落水，在道德上男友有救助义务，但法律没有给男友规定这项义务。因此，狗蛋不救助母亲，构成不作为犯罪。如果妻子与母亲都落水，丈夫（儿子）对妻子和母亲都有救助义务。如果丈夫救了妻子，未救母亲，则对母亲属于不作为，然后在排除违法事由阶段，根据“义务冲突”这种排除违法事由排除不作为的违法性。\footnote{所谓义务冲突，是指行为人负有两项义务，履行其中一项义务，必然无法履行另一项义务；未履行另一项义务，不具有违法性。例如，家庭医生甲发现男主人感染了艾滋病。甲既有为男主人保密的义务，又有告知女主人的义务。甲选择告知女主人，那么未履行保密义务，不具有违法性。}（2015年第52题，2020年试题）

[巩固练习］题1，甲带邻居小孩出去玩，小孩掉入粪池。甲虽然不是小孩的父亲，但属于临时的监护人，有保护义务。（2010年第52题）

题2，法官作出甲乙离婚的一审判决书。甲乙走出法院，乙遭遇车祸。甲有救助义务，因为判决尚未生效。（2024年试题）

[注意］多个作为义务人之间的先后关系。多数说：若第一位的作为义务人不救助法益甚至侵害法益，则不能期待第二位的作为义务人实施救助。

例1，岳母受伤，妻子不救，丈夫也不救。妻子是第一作为义务人，丈夫是第二作为义务人。丈夫不构成不作为犯。

例2，妻子殴打自已母亲，丈夫不阻止，不构成不作为犯。（2020年试题）

（2）职务、业务产生的特定关系

例如，警察对犯罪行为中的被害人、消防队员对火灾中的被害人、医生对病人、游泳池的救生员对游泳者，都负有保护义务。当然，警察、消防队员的这种职务也有法律规定，因此也属于法律规范产生的保护义务。

例1，民警甲在辖区内酒吧喝酒，发现有人在强奸乙女。甲认为自己已下班，遂未制止。甲构成不作为犯罪，因为《警察法》规定，警察在非工作时间，遇有职责范围内的紧急情况，应当履行职责。这份职责，不受是否穿警服的影响。（2023年试题）

例2，甲在杀乙，卫生局局长丙路过。丙虽然是公务员，但不是警察，在职务上没有保护乙的义务。（2006年第4题）

（3）合同约定产生的特定关系。例如，顾客甲在试衣间试衣服，将手机放在柜台，让店员乙照看，乙答应。乙看到丙偷手机，不阻止。基于约定，乙有保护义务。乙构成盗窃罪的不作为的帮助犯。（2025年试题）

（4）\CJKunderwave{基于自愿救助产生的特定关系}。这是指某项法益处于危险境地时，行为人自愿救助，使法益的保护依赖于行为人时（形成依赖关系），行为人就有继续保护的义务。总结产生作为义务的条件：第一，自愿救助。第二，法益对象对该行为产生依赖关系。

例1，甲对路边弃婴视而不见，不构成犯罪，若捡回家，则有继续救助的义务。虽然

% Source PDF physical page 50; see audit/source_coverage.json.
捡回家是降低危险的行为，但是婴儿对甲形成了依赖关系。（2013年第51题）

例2，甲请钟点工乙照顾婴儿，到了工作结束时间，甲还未回家。乙看到婴儿有危险，但认为到了合同约定的下班时间，便离去。婴儿受伤。虽然在合同约定上，钟点工乙没有继续救助义务，但是，乙基于平等自愿而照顾婴儿，那么当婴儿有危险，婴儿的父母还未回家，此时婴儿对乙产生依赖关系，乙便有继续照顾婴儿的义务。

[巩固练习］题1，甲给井底的乙扔绳子，绳子到了井底，甲尚未往上拉时，发现乙是仇人，便不拉绳子。由于甲尚未拉绳子，乙对甲尚未形成依赖关系，因此甲没有继续救的义务。如果甲将乙拉到半空时，便有继续往上拉的义务。所谓救人救到底，送佛送到西。如果此时甲放掉绳子，乙摔死，则甲同时构成作为与不作为，二者竞合，最终认定为作为。（2014年第5题）

题2，甲发现乙溺水，扔下救生圈。在乙尚未抓到救生圈时，甲发现乙是自已的仇人，于是捞走救生圈，乙溺亡。乙尚未抓到救生圈，尚未形成依赖关系，所以甲没有救助义务。（2024年试题）

题3，甲请乙到酒吧喝酒。乙自己喝醉。第一，甲直接离开。乙自己回家，因为醉酒找不到家，冻死在路上。甲没有救助义务，因为甲请乙喝酒的行为属于生活行为，而不是给乙的法益创设危险的先行行为。该死亡结果与甲没有因果关系（2025年试题）。此时，根据后文的特定领域，酒吧有保护义务。第二，甲实施了救助行为，同时乙依赖甲的救助行为，则甲有继续救助的义务。例如，甲将乙扶上车，开车送乙回家，此时便有继续送回家的义务。\footnote{参见张明楷：《刑法的私垫（之二）》，北京大学出版社2017年版，第40页。}

题4，甲乙丙相约探险、登山（比较危险的山），意味着愿意相互照顾，形成\CJKunderwave{危险共}\CJKunderwave{同体}，这种特定关系会产生保护义务。丙遇到危险时，甲乙有救助义务。（2023年试题）

题5，甲乙相约打羽毛球。球场上，甲突发心脏病。乙没有刑法上的救助义务，因为没有形成危险共同体。（2024年试题）

2.\CJKunderwave{特定领域}

基于特定领域对被害人产生保护义务，需要具备两个条件：第一，行为人是特定领域的管理者，被害人不是特定领域的管理者。二者是主人与客人的关系。第二，行为人对特定领域内的危险具有控制力，被害人对行为人形成依赖关系。

例1，甲学生和乙学生共同居住在某宿舍。某日，甲发现屋内着火，没有救助乙，自已逃出宿舍。甲没有刑法上的救助义务。虽然甲是特定领域的管理者，但是乙也是特定领域的管理者。（2023年试题）

例2，案。第一，地点在一方家里。卖淫女与客在卖淫女家里发生性交易。过程中客心脏病发作，卖淫女有救助义务。卖淫女心脏病发作，客没有救助义务。如果地点在客家里，则分析原理同上。第二，地点在酒店，看是谁开的房。谁开的房，则视为在谁家里，分析原理同上。此时，酒店管理者基于职务、业务也有救助义务。第三，地点在二人共同管理的房间，则另一方没有救助义务。（2023年试题）

例3，吸毒者甲乙经常在乙家吸毒。有次，乙吸毒后，口吐白沫。甲没有救助义务。如果甲吸毒后，口吐白沫，则乙有救助义务。（2023年试题）

% Source PDF physical page 51; see audit/source_coverage.json.
[巩固练习］题1，出租车司机看到车上男乘客在强奸女乘客，司机负有阻止义务。这是因为，出租车内是个封闭空间，司机对该空间内的危险的控制力较强，女乘客的获救依赖于司机。司机如果不阻止，就构成强奸罪的不作为形式的帮助犯。

题2（2023年试题），公交车司机在驾车途中，看到小偷盗窃乘客的财物。司机没有刑法上的阻止义务。这是因为，公交车内是个具有开放性、流动性的公共空间，司机对该空间内的危险的控制力较弱，乘客具有其他获救机会，对司机没有形成依赖关系。

题3-1，甲开车肇事，导致乙受伤。甲拦住出租车后，将乙搬入出租车内准备送往医院，途中甲借故下车逃离。乙处在出租车司机丙独立支配的车内，丙负有救助义务。当然，这并不意味着要求丙对乙负完全责任，只是要求丙此时不应将乙遗弃。

题3-2，甲开车不慎撞上出租车，导致出租车上的乘客乙受伤。甲与出租车司机丙均有救助乙的义务。甲如果逃逸，则丙需要救助乙。甲的作为义务来源于先行行为。丙的作为义务来源于特定领域和业务关系。

题4，甲家封闭院落闯入一个危重病人乙，甲有救助义务。履行该义务，不意味着要对乙负完全责任，只需要将乙送医院，或移至其他人能够发现的场所。意思是，乙的小命在甲手里，甲不能看着乙死。

\InsertDiagram{figures/manual_p051_1.tex}

\subsection*{二、具有作为可能性 （能为）}

刑法给行为人增设某项作为义务是有目的的，要追求一种效果、功能，也即行为人能够履行作为义务，以避免法益侵害结果发生。这便是作为义务的实效性。作为义务要具有实效性的第一项条件是行为人要具有履行能力（作为可能性）。“超出能力，则义务失效”。未履行失效的义务，不具有违法性，因此不构成不作为犯罪。

例如，孩子落水，父亲不会游泳，无法救孩子。虽然父亲负有作为义务，但是缺乏履行能力。这意味着，此时作为义务对父亲失去实效性。父亲未履行失去实效的作为义务，不具有违法性，因此不构成不作为犯罪。

\subsection*{三、具有结果避免可能性（不为）}

作为义务要具有实效性，第二项条件是结果避免可能性，也即如果行为人履行了作为义务，那么具有结果避免发生的可能性。如果行为人努力履行义务，结果仍然会发生，那么意味着该作为义务此时失去实效性。没有履行失效的义务，不具有违法性，因此不构成不作为犯罪。

% Source PDF physical page 52; see audit/source_coverage.json.
例1，甲在车间操作不当，导致机器压伤乙，致乙濒临死亡。甲不救乙，乙很快死亡。即使甲努力施救，也回天无力。甲不构成不作为犯罪，只构成过失致人死亡罪。（2010年第52题）

例2，交通肇事罪有个法定刑升格条件“因逃逸致人死亡”。如果伤者被撞成头盖骨破裂，濒临死亡，即使立即送最近的医院抢救也无法挽救生命，这种情况下肇事者逃逸，伤者死亡的，肇事者不属于“因逃逸致人死亡”。

\subsection*{四、主观要件}

前述三项条件是不作为犯的客观条件。最终成立不作为犯还需要主观上有故意或过失。就故意的不作为犯而言：

\subsubsection*{（一）事实认识错误}

行为人对产生作为义务的基础事实存在认识错误，属于关于事实的认识错误，会排除行为人的犯罪故意。

例如，甲路过河边，看到有人落水，误以为是无关的陌生人落水，没有施救，实际上是自己的小孩落水，小孩溺亡。甲没有认识到是自已的小孩（作为义务的基础事实），不是故意犯罪，但如果有认识的可能性，存在过失，则甲成立过失的不作为犯罪，也即不作为的过失致人死亡罪（以符合不作为犯的客观要件为前提）。

\subsubsection*{（二）法律认识错误}

行为人对产生作为义务的基础事实没有认识错误，但对是否产生刑法上的作为义务存在认识错误，属于关于法律的认识错误，不会排除行为人的犯罪故意，不能免责。

例如，甲路过河边，看到自己的小孩（养子）落水，误以为自己没有救助义务，没有施救，小孩溺亡。甲的认识错误属于法律认识错误，误以为自已没有刑法上的义务，误以为不救助不构成犯罪。然而，不知法者不免责。甲成立故意的不作为犯罪（遗弃罪或不作为的故意杀人罪）。一般人只要认识到是自已的孩子、妻子、父母，就会认识到有刑法上的救助义务。甲不可能认识不到这一点，甲具有认识到这一点的可能性。（2015年第52题）

[提示1」等价性问题。由于给不作为定罪，有点“无中生有”，所以，为了不违反罪刑法定原则，要给某个不作为定罪，要求该不作为能够达到与作为犯罪相提并论的程度，能够相同评价（等价性）。例如，母亲故意不喂养婴儿，持续时间是一个小时。此时，程度尚未达到犯罪程度。如果持续三天，饿死婴儿，那么在等价性上，达到了作为的故意杀人罪的程度，可以定不作为的故意杀人罪。

[提示2］故意的不作为犯也有既遂、未遂形态。例如，甲将自己的婴儿饿了三天，婴儿快要死了，邻居发现，赶紧抢救了婴儿。甲构成不作为的故意杀人罪的未遂。

% Source PDF physical page 53; see audit/source_coverage.json.
\OriginalPanel{典型真题}

关于不作为犯罪，下列哪些选项是正确的？（2015年·卷二·52题）\footnote{[答案]B项，甲属于法律认识错误，不排除犯罪故意，也不能免责，成立故意的不作为犯罪。C项，在法律上，甲对母亲有救助义务，但对女友没有救助义务。在客观要件阶段，甲不救助母亲构成不作为。如果存在排除违法事由，则可以排除甲的这种不作为的违法性；如果不存在，则甲最终构成不作为犯罪。应注意的是，如果将题干改为“母亲与老婆”，则性质不同。在法律上丈夫对老婆也有救助义务。此时丈夫不救母亲，在客观要件阶段先构成不作为，然后根据义务冲突这种排除违法事由排除不作为的违法性，最终不构成不作为犯罪。D项，甲对丙具有作为义务也即制止的义务，作为义务来自甲的先行行为，也即投毒杀乙的行为。甲就在现场，具有履行义务的可能性却不履行，构成不作为犯罪。本题答案：ACD。}

A.儿童在公共游泳池溺水时，其父甲、救生员乙均故意不救助。甲、乙均成立不作为犯罪

B.在离婚诉讼期间，丈夫甲误认为自己无义务救助落水的妻子，致妻子溺水身亡的，成立过失的不作为犯罪

C.甲在火灾之际，能救出母亲，但为救出女友而未救出母亲。如无排除犯罪的事由，甲构成不作为犯罪

D.甲向乙的咖啡投毒，看到乙喝了几口后将咖啡递给丙，因担心罪行败露，甲未阻止丙喝咖啡，导致乙、丙均死亡。甲对乙是作为犯罪，对丙是不作为犯罪

